\documentclass[10pt,a4paper]{amsart}
\usepackage[margin=19mm]{geometry}
\usepackage{amsmath,amssymb,mathtools}
\usepackage[scr=rsfs,cal=euler]{mathalfa}
\usepackage{xfrac}
\usepackage[unicode,hidelinks,bookmarks=false]{hyperref}
\newcommand{\ie}{i.e.\ }
\newcommand{\cf}{cf.\ }
\newcommand{\resp}{respectively\ }
\newcommand{\Subsection}{\subsection}
\providecommand{\nobreakdash}{\nobreak\hbox{-}}
\setlength{\emergencystretch}{2em}
\multlinegap0pt
\usepackage{amssymb}
\newtheorem{theorem}{Theorem}
\def\IR{{\mathbb R}}
\def\IV{{\mathbb V}}
\def\dis{\displaystyle}
\def\f{\footnotesize}
\def\n{\noindent}
\def\wh{\widehat}
\def\wt{\widetilde}
\title{An effective energy-enstrophy diffusion process with a condensation bound}
\author{Alain-Sol Sznitman, Klaus Widmayer}
\date{Source: arXiv:2602.15810v2 (CC BY 4.0)}
\begin{document}
\maketitle

\noindent\emph{Source and licence.} An effective energy-enstrophy diffusion process with a condensation bound. Version arXiv:2602.15810v2 (CC BY 4.0), distributed by arXiv under the Creative Commons Attribution 4.0 International licence, http://creativecommons.org/licenses/by/4.0/, as recorded in the arXiv licence metadata for this exact version and shown on the abstract page https://arxiv.org/abs/2602.15810v2. Full licence notice: \texttt{licenses/arxiv-2602.15810-CC-BY-4.0.txt}.\par\smallskip

\noindent\textit{Editorial scope.} Counted unit: the article's theorem theo2.3 (source0.tex lines 398--404). The article's complete original proof of it is delivered with it (source0.tex lines 406--491, one \texttt{proof} environment of 86 lines). No mathematics is written, repaired or re-derived: the statement and the proof are the article's own lines, in the article's own notation, and no retained line is substituted or re-flowed.  Retained uncounted as the source-local context the proof needs: lines 253--255; lines 265--267; lines 269--271; lines 340--344; lines 1023--1030; lines 1375--1383; lines 1385--1391.  Nothing was omitted from any retained span, so the package carries no omission record.  No retained line contains a citation, so the wrapper carries no bibliography and no bibliography entry.  No retained line contains a figure environment, an image-inclusion command or drawing code, so no image is delivered and the package declares no asset dependency.  Editorial formatting only, in addition: an AMS \texttt{amsart} preamble that restates the article's own declarations and macros used by the retained lines, namely the environments \texttt{theorem} on separate counters, together with the standard packages those lines need.  The retained environments print in this document as Theorem 1 in order of appearance, whereas the article prints the same items with the numbering of its own full document.  The complete native source of the article as distributed in the arXiv e-print for this exact version is bundled unchanged as \texttt{sources/194829/source0.tex}.
\begin{equation}\label{1.1}
N = 2 n, \; \mbox{with $n \ge 4$}
\end{equation}

\begin{equation}\label{1.4}
\mbox{$\mu_i = \lambda_{2i}$, for $1 \le i \le n$ (and $n$ as in (\ref{1.1}))}.
\end{equation}

\begin{equation}\label{1.5}
a > 0,
\end{equation}

\begin{align}
& q_{2i} = q_{2i-1}, \; \mbox{for $1 \le i \le n$ (and we write $\wh{q}_i = q_{2i} + q_{2i - 1})$}, \label{2.2}
\\ 
& u = \sum_1^N \,q_\ell(u,v), \, v = \sum_1^N \; \dis\frac{1}{\lambda_\ell} \;q_\ell (u,v), \; \mbox{for all $(u,v) \in C$}. \label{2.3}
\end{align}

\begin{align}
& \Sigma = \{3, \dots, n\}, \label{A.4}
\\[1ex]
&\IV_{u,v} = \Big\{(s_3, \dots, s_n) \in \IR^\Sigma_+; \; u - \mu_1 v \ge \Sigma^n_3 \, \Big(1- \mbox{\f $\dis\frac{\mu_1}{\mu_i}$}\Big)\,s_i \; \mbox{and}
 \label{A.5}
\\
&\qquad \qquad \qquad\qquad  \qquad \quad u - \mu_2 v \le \Sigma^n_3 \, \Big(1-\mbox{\f $\dis\frac{\mu_2}{\mu_i}$}\Big)\,s_i\Big\}, \nonumber
\end{align}

\begin{equation}\label{B.16}
\begin{array}{l}
\mbox{the vector $(\wh{q}_i(u,v))_{i \in \Sigma}$ is the expectation of $(S_i)_{i \in \Sigma}$ under the conditional}
\\
\mbox{density of $(S_i)_{i \in \Sigma}$ given $U = u$ and $V = v$, and is also the barycenter of $\IV_{u,v}$}
\\
\mbox{(and hence does not depend on $a$ in (\ref{1.5}))}.
\end{array}
\end{equation}

\begin{equation}\label{B.17}
\begin{split}
\mu_2 (1-1/\mu_2) \, \wh{q}_1(u,v) = &\; - u + \mu_2 v + \Sigma^n_3 \,(1- \mu_2/\mu_i) \, \wh{q}_i(u,v)
\\
(1-1/\mu_2) \, \wh{q}_2(u,v) = &\; u - \mu_1 v - \Sigma^n_3 \,(1- \mu_1/\mu_i) \, \wh{q}_i(u,v),
\end{split}
\end{equation}

\begin{theorem}\label{theo2.3} {\it (Monotonicity and upper bound)}
\begin{align}
& \mbox{$(1 - \mu_i^{-1}) \; \wh{q}_i (u,v)$ is non-decreasing in $i$ for $2 \le i \le n$ and $(u,v) \in C$}. \label{2.11}
\\[1ex]
& \wh{q}_i \, (u,v) \le [(n - i + 1) (1 - \mu_i^{-1})]^{-1} (u-v), \;\mbox{for $2 \le i \le n$ and $(u,v) \in C$}. \label{2.12}
\end{align}
\end{theorem}

\begin{proof}
Note that by (\ref{2.3}) one has
\begin{equation}\label{2.13}
\dis\sum^n_2 \;(1 - \mu^{-1}_i) \; \wh{q}_i \,(u,v) = u-v, \; \mbox{for all $(u,v) \in C$}.
\end{equation}
All the summands in the left member above are non-negative and (\ref{2.12}) is an immediate consequence of the monotonicity property (\ref{2.11}). To complete the proof of Theorem \ref{theo2.3}, we thus only need to show (\ref{2.11}). Due to the continuity of the $\wh{q}_i$, it suffices to prove (\ref{2.11}) when $(u,v) \in \; \stackrel{_\circ}{C}$.

\medskip
To this end, we introduce $\wt{\Sigma} = \{2,\dots,n\}$ (we recall that $\Sigma = \{3,\dots, n\}$, see (\ref{A.4}) and that $\mu_1 = 1$, see (\ref{1.4})) and define the affine isomorphism $h$ between $\IR^\Sigma$ and the affine hyperplane $\wt{H} = \{\wt{\sigma} = (\wt{s}_i)_{2 \le i \le n} \in \IR^{\wt{\Sigma}}$; $\Sigma_2^n \, \wt{s}_i = u-v\}$:
\begin{equation}\label{2.14}
\begin{array}{l}
h(\sigma) = \wt{\sigma} \; \mbox{with} \; \sigma = (s_i)_{3 \le i \le n}, \; \wt{\sigma} = (\wt{s}_i)_{2 \le  i \le n}, \; \mbox{and}
\\[1ex]
\wt{s}_i = (1 - \mu^{-1}_i) \; s_i, \; \mbox{for} \; 3 \le i \le n, \; \wt{s}_2 = u-v - \sum^n_3 (1- \mu_i^{-1}) \, s_i .
\end{array}
\end{equation}
Then, $h$ maps $\IV_{u,v} = \{(s_3,\dots, s_n) \in \IR^\Sigma_+$; $u - \mu_1 v \ge \Sigma^n_3 \,(1 - \frac{\mu_1}{\mu_i})\, s_i$ and $u - \mu_2 v \le \Sigma^n_3 \,(1 - \frac{\mu_2}{\mu_i}) \,s_i\}$, see (\ref{A.5}), one-to-one onto
\begin{equation}\label{2.15}
\begin{array}{l}
\wt{\IV}_{u,v} = \{ \wt{\sigma} = (\wt{s}_i)_{2 \le i \le n} \in \IR^{\wt{\Sigma}}_+; \; \Sigma_2^n \, \wt{s}_i = u-v \;\,\mbox{and} \; \,\Sigma_2^n \, \gamma_i \,\wt{s}_i \ge u - \mu_2 v\} \subseteq \wt{H},
\\[1ex]
\mbox{with} \; \gamma_i = \mbox{\f $\dis\frac{1 - \mu_2/\mu_i}{1 - \mu_1/\mu_i}$} = 1 -  \mbox{\f $\dis\frac{\mu_2-\mu_1}{ \mu_i - \mu_1}$} , \; \mbox{for $2 \le i \le n$ (and we recall that $\mu_1 = 1)$}.
\end{array}
\end{equation}
Note that the $\gamma_i$ are non-decreasing in $i$ and that $\gamma_2 = 0$. Further, $h$ maps the Lebesgue measure on $\IR^\Sigma$ into a constant multiple of the surface measure on the hyperplane $\wt{H}$, so that $h$ maps the barycenter of $\IV_{u,v}$ on the barycenter of $\wt{\IV}_{u,v}$, which we denote by $\wt{q}(u,v) = (\wt{q}_i (u,v))_{2 \le i \le n}$. We thus see by (\ref{B.16}) that
\begin{equation} \label{2.16}
 h\big(\wh{q}(u,v)\big) = \wt{q}(u,v), \; \mbox{with the notation $\wh{q} (u,v) = \big(\wh{q}_i (u,v)\big)_{3 \le i \le n}$, i.e.}
 \end{equation}
\begin{equation}\label{2.17}
\begin{split}
 \wt{q}_i (u,v)  &= (1 - \mu^{-1}_i) \, \wh{q}_i (u,v), \; \mbox{for $3 \le i \le n$},
\\
 \wt{q}_2(u,v) &= u-v - \Sigma^n_3 \,(1 - \mu_i^{-1})\, \wh{q}_i(u,v) \stackrel{(\ref{B.17})}{=} (1 - \mu^{-1}_2) \; \wh{q}_2 (u,v).  
 \end{split}
\end{equation}

\medskip\n
The claim (\ref{2.11}) thus amounts to showing that
\begin{equation}\label{2.18}
\mbox{$\wt{q}_i (u,v), 2 \le i \le n$, is non-decreasing in $i$}.
\end{equation}

\n
For this purpose we consider $i_0 \in \{2, \dots, n-1\}$ and the map $S$ swapping the $i_0$ coordinate with the $i_0 + 1$ coordinate on $\IR^{\wt{\Sigma}}$, namely, such that for $\wt{\sigma} = (\wt{s}_i)_{2 \le i \le n}$
\begin{equation}\label{2.19}
\begin{split}
S(\wt{\sigma})_i & = \wt{s}_i, \; \mbox{when $i \in \wt{\Sigma}\; \backslash\,\{i_0, i_0 + 1\}$},
\\
& = \wt{s}_{i_0 + 1}, \; \mbox{when $i = i_0$},
\\
& = \wt{s}_{i_0}, \; \mbox{when $i = i_0 + 1$}.
\end{split}
\end{equation}
We also introduce (see (\ref{2.15}) for notation)
\begin{equation}\label{2.20}
\wt{\IV}_{i_0, +} = \{\wt{\sigma} \in \wt{\IV}_{u,v}; \; \wt{s}_{i_0} < \wt{s}_{i_0 + 1}\}, \; \wt{\IV}_{i_0,-} = \{\wt{\sigma} \in \wt{\IV}_{u,v}; \; \wt{s}_{i_0 + 1} < \wt{s}_{i_0}\}.
\end{equation}
We then note that
\begin{equation}\label{2.21}
\begin{array}{l}
\mbox{$\wt{\IV}_{i_0,+}$, $\wt{\IV}_{i_0,-}$ are pairwise disjoint convex polytopes contained in $\wt{\IV}_{u,v}$}, 
\\
\mbox{and their union differs from $\wt{\IV}_{u,v}$ by a negligible set}.
\end{array}
\end{equation}
The key observation is now that
\begin{equation}\label{2.22}
\mbox{$S$ maps $\wt{\IV}_{i_0, -}$ into $\wt{\IV}_{i_0, +}$},
\end{equation}
because, see (\ref{2.15}) for notation,
\begin{equation}\label{2.23}
\gamma_{i_0} \, \wt{s}_{i_0} + \gamma_{i_0 + 1} \, \wt{s}_{i_0 + 1} < \gamma_{i_0} \, \wt{s}_{i_0 + 1} + \gamma_{i_0 + 1} \, \wt{s}_{i_0} \; \mbox{when} \; \wt{s}_{i_0 + 1} < \wt{s}_{i_0},
\end{equation}

\n
since  their difference equals $(\gamma_{i_0 + 1} - \gamma_{i_0}) (\wt{s}_{i_0} - \wt{s}_{i_0 + 1}) > 0$, so that $S(\wt{\IV}_{i_0, -}) \subseteq \wt{\IV}_{u,v}$, and (\ref{2.22}) readily follows.

\medskip
The map $S$ preserves the surface measure on $\wt{H}$, and we thus see that
\begin{equation}\label{2.24}
\mbox{the barycenter $\wt{q}(u,v)$ of $\wt{\IV}_{u,v}$ belongs to the closure of $\wt{\IV}_{i_0, +}$}.
\end{equation}

\n
This shows that $\wt{q}_{i_0}(u,v) \le \wt{q}_{i_0 + 1}(u,v)$ for all $i_0$ in $\{2,\dots,n\}$ and $(u,v)$ in $\stackrel{_\circ}{C}$. The claim (\ref{2.18}) follows, and as already observed this completes the proof of Theorem \ref{theo2.3}.
\end{proof}

\end{document}
